\chapter{Outiller la productivité et standardiser l'environnement de développement}
\label{chap:outils}
% Chapitre-chapeau du volet 1 : Grafana (t1), compilation (t2), DSDK (t6).

Ce premier chapitre regroupe trois missions qui partagent une même nature :
construire ou reprendre en main l'\textbf{outillage} qui conditionne le
travail quotidien des développeurs. Relues à travers la typologie du
chapitre~\ref{chap:veille}, elles s'attaquent principalement à la
\textbf{dette de build et de dépendances} et à la \textbf{dette de
connaissance} : des procédures manuelles, non reproductibles ou détenues par
une seule personne, qui renchérissent chaque intervention sur le logiciel.

L'intérêt de \textbf{standardiser l'environnement de développement} est
double. C'est d'abord un enjeu de \textbf{productivité} : tout ce qu'un
développeur doit refaire à la main (répliquer une configuration de
supervision, subir une compilation qui sature sa machine, pousser le code
vers une machine distante pour le construire) est un coût payé à chaque
itération, par chaque membre de l'équipe ; l'outiller, c'est le payer une
fois. Ce coût se chiffre : à l'échelle d'une journée de développeur, qui
recompile et redéploie des dizaines de fois, une à deux minutes gagnées par
cycle se comptent presque en \textbf{heures} récupérées, soit du temps passé à
produire plutôt qu'à attendre. C'est ensuite un enjeu de \textbf{capitalisation} : un environnement
standardisé transforme un savoir-faire individuel en savoir
\emph{exécutable} (outil, configuration versionnée, image de conteneur) qui
survit au départ des personnes, là où la procédure manuelle disparaît avec
son détenteur. Enfin, un socle commun à toute l'équipe supprime les écarts
entre postes (« ça compile chez moi ») qui transforment chaque anomalie en
enquête.

Trois missions actionnent ce levier : un outil de duplication de la
configuration Grafana (section~\ref{chap:t1}) et la maîtrise du coût de
compilation du \emph{Valo Engine} (section~\ref{chap:t2}), toutes deux chez
Free Mobile, puis, chez Beemo, la rétro-ingénierie du DSDK pour rendre le
\emph{build} de l'agent VSX local, reproductible et documenté
(section~\ref{chap:t6}).
